\section{Оценка компонентов: дизайн и оцениваемые эффекты}
\label{sec:heldout-methods}

\subsection{Факторный эксперимент без компрессии}
Основной эксперимент оценивает ядро Hybrid-INoT и его контроли: независимо варьируются число вызовов модели (один или три) и обозначения этапов (нейтральные или ролевые). Таблица~\ref{tab:design} определяет четыре факторных условия и прямой решатель. Во всех структурированных условиях предписаны одинаковые операции: анализ требований и граничных случаев, реализация по плану, затем проверка и исправление относительно требований. Ролевое воздействие заменяет обозначения «stage 1», «stage 2», «stage 3» на «planner», «implementer», «reviewer»; предложения с описанием операций остаются неизменными. При одном вызове все три инструкции предшествуют задаче. При трёх вызовах каждый получает инструкцию своего этапа, полный текст задачи, весь предоставленный контекст и все полные предыдущие видимые ответы. Итоговый формат одинаков: ровно один блок с полной программой Python. Прямой решатель получает только инструкцию реализации и требование к итоговому формату.

\begin{table}[htbp]\centering
\caption{Реализованные условия. Фактор числа вызовов включает передачу промежуточных ответов между новыми вызовами.}\label{tab:design}
\begin{tabular}{llll}\toprule
Код & Вызовы & Обозначения & Операции\\\midrule
D & 1 & Нет & Прямая реализация\\
SN & 1 & Нейтральные & План, реализация, проверка\\
SR & 1 & Ролевые & План, реализация, проверка\\
MN & 3 & Нейтральные & План, реализация, проверка\\
MR & 3 & Ролевые & План, реализация, проверка\\\bottomrule
\end{tabular}\end{table}

Компрессия отсутствует во всех условиях. Такая постановка заменяет прежнее сравнение трёх вызовов без сжатия с одним сжатым вызовом, но не оценивает эффект прежнего компрессора. Контраст обозначений является воздействием на инструкцию, а не свидетельством независимых внутренних агентов. Контраст числа вызовов включает поток промежуточной информации и повторный инференс, а не только механическое изменение счётчика вызовов. Непроверенные маршрутизатор и дополнительная модель проверки исключены из заявленного метода.

Пусть $Y_t(a)$ показывает, прошла ли единственная итоговая программа для задачи $t$ в условии $a$ исходные тесты бенчмарка в проверенном окружении. Четыре заранее заданных парных эффекта качества имеют вид
\begin{align}
\Delta_{R,1}&=\E_t[Y_t(\mathrm{SR})-Y_t(\mathrm{SN})], &
\Delta_{R,3}&=\E_t[Y_t(\mathrm{MR})-Y_t(\mathrm{MN})],\\
\Delta_{C,N}&=\E_t[Y_t(\mathrm{MN})-Y_t(\mathrm{SN})], &
\Delta_{C,R}&=\E_t[Y_t(\mathrm{MR})-Y_t(\mathrm{SR})].
\end{align}
Каждая нулевая гипотеза задаёт нулевую разность качества. Описательное взаимодействие равно $\Delta_{R,3}-\Delta_{R,1}$. Ожидания относятся к выборке задач при реализованном протоколе и стохастической генерации, а не к свойству, независимому от поставщика модели. Оценки по доступным данным требуют полных пар; при пропусках без дополнительных предположений они характеризуют именно наблюдаемую подвыборку.

\subsection{Распределение задач и отдельный этап разработки}
Источник задач --- BigCodeBench v0.1.4~\cite{zhuo2024bigcodebench}. Основная выборка содержит 200 новых идентификаторов из ранее зафиксированной случайной перестановки резервного пула с заданным начальным состоянием генератора. Взяты первые 200 допустимых идентификаторов в этом случайном порядке после исключения восьми примеров карточки набора данных (/0--/7), случайно просмотренных при проверке лицензии. Это не первые строки датасета и не отбор по сложности. Идентификаторы, точные входы модели и данные оценщика хешируются до контролей; неудачный контроль не приводит к замене задачи. Сорок задач разработки не пересекаются с этими 200. Этап разработки включает два обозначенных повтора и 400 программ, основная выборка --- один новый повтор и 1\,000 назначений. Последние дают не более 200 независимых единиц выборки, а не 1\,000 независимых задач. Родственные задачи могут дополнительно уменьшать эффективную независимость. Публичность бенчмарка оставляет открытым вопрос о наличии задач в обучающих данных.

Назначение означает запланированную генерацию для заданной задачи, условия и повтора; начатое назначение считается попыткой. Каждой основной задаче назначены все пять условий. Четыре непересекающиеся группы по остатку индекса от деления на четыре содержат по 50 задач. В каждой группе зафиксированный псевдослучайный порядок определяет блоки задач и порядок условий; группу обрабатывают два работника, одновременно действует не более восьми экспериментальных процессов CLI. Полная матрица и хеши исходного кода закоммичены до отправки запросов. Модель не получает тесты бенчмарка, эталонные решения или результаты тестирования во время генерации. Метки повторов обозначают новые наблюдения по протоколу, а не случайные состояния модели у поставщика.

\subsection{Генерация и полные доступные трассы}
Используются запрошенный псевдоним модели \texttt{gpt-5.4-mini}, уровень рассуждения medium и официальный Codex CLI 0.153.4 с аутентификацией подписки ChatGPT~\cite{openai54mini2026,codexexec2026}. Опубликованный тариф позволяет проверить денежную оценку ресурсов; выбранная конфигурация является экономичным доступным вариантом для кода, но глобальный минимум цены не доказан. CLI изолирован от пользовательской конфигурации и инструкций репозитория. Инструменты, браузер, делегирование и выполнение тестов отключены. Каждый вызов создаёт новую временную сессию в пустом рабочем каталоге. Допустимая схема событий требует ровно одного полного ответа и корректной записи использования; неожиданный инструмент, повтор или событие сжатия нарушают заданное текстовое условие.

Для каждого отправленного вызова сохраняются точный промпт, аргументы, поток событий, стандартный поток ошибок и конечное состояние. Для завершённых записей дополнительно сохраняются итоговый текст, компоненты числа токенов и хеши исходных файлов. Входы последующих вызовов реконструируются и побайтно сверяются с полными предыдущими доступными ответами. Доступность скрытых рассуждений поставщика не предполагается. Псевдоним модели не фиксирует неизменные веса; проверенное управление температурой и случайным состоянием модели отсутствует. На вызов установлен предел 180 секунд. Вход свыше 65\,536 байт UTF-8 отклоняется, а не сокращается. Единого жёсткого ограничения выходных токенов нет; длина ответа и потребление ресурсов являются результатами воздействия.

Исходное правило останавливает группу после первой неудачной попытки. Поэтому сетевые обрывы и превышения времени оставляют другие назначения неотправленными. Отдельная закоммиченная поправка, записанная до оценки или просмотра качества основной выборки, разрешает продолжить только такие неотправленные назначения. Любое назначение хотя бы с одним начатым этапом навсегда исключается из продолжения. Точный список оставшихся назначений и завершённый архив первой серии фиксируются до отправки; сохраняются промпты, настройки, ограничения времени и исходный порядок внутри группы. Явный сетевой обрыв или превышение времени завершает попытку без повтора; ошибки квоты, аутентификации и нераспознанные ошибки останавливают дальнейшую отправку. Первые попытки, продолжения и пропуски остаются различимыми. Таким образом, анализ использует заранее заданные контрасты при изменённом протоколе исполнения; это не неизменённый пререгистрированный запуск.

\subsection{Штатная оценка и контроль оцениваемости}
Неизменённый CLI BigCodeBench на коммите \texttt{09dd993f46c3} выполняет исходные тесты задач в режиме \texttt{instruct full}. Каждый запуск оценщика использует контейнер без сети, от непривилегированного пользователя, с основной файловой системой только для чтения, пределом памяти 3 ГБ, двумя CPU и двумя работниками оценщика. Подготовленные данные, точное дерево исходников, идентификатор образа, версии пакетов, автономные ресурсы, команды и штатные отчёты архивируются. Строгое извлечение выбирает единственный итоговый блок Python; нарушение формата даёт пустую наблюдаемую программу и отдельный признак ошибки формата, без выбора альтернативного ответа или исправления.

До генерации модели исходные эталонные и заведомо неверные программы проверены на всех 200 задачах. Изменения зависимостей ограничены этапом контролей, сохраняют предыдущие неудачные запуски и вызывают повторные контроли всей матрицы. Итоговое окружение пропускает 193 эталона и отклоняет все 200 неверных программ. Поэтому семь задач недоступны для оценки качества во всех условиях; их назначения и ресурсы генерации сохраняются в учёте. Итоговый образ дополняет окружение разработки фиксированными зависимостями и автономными ресурсами NLTK. Штатная диагностика эталонов содержит ошибки DNS (/101, /590), отсутствие TensorFlow (/418), неподдерживаемый аргумент кодировщика (/686), проверки моментов вырожденной выборки (/276) и проверки архива/журнала (/14, /1028). Последующая проверка без изменения образа подтверждает отсутствие внешних команд, используемых последними двумя задачами; связь этих отказов с отсутствием команд является диагностическим выводом по исходникам. Позднейшая диагностика не меняет зафиксированный контроль оцениваемости. Контроли устанавливают оцениваемость и различение заведомо неверного решения, но не полную семантическую согласованность инструкции и теста.

Оценщику передаются только реально полученные программы. Каждый результат связывается с точной архивной программой, исходным идентификатором задачи и проверенным окружением; незавершённые запуски оценщика и противоречивые отчёты не могут давать наблюдения качества. Исходный штатный статус хранится рядом с отдельным аналитическим статусом с учётом контроля оцениваемости. Для отсутствующей генерации не создаётся фиктивный отказ тестов. Одинаковый контроль применяется ко всем условиям и поисковой репликации INoT.

\subsection{Статистический анализ, пропуски и оценка ресурсов}
Для каждого из четырёх запланированных контрастов качества приводятся число полных пар, оба числа несогласованных исходов, средняя парная разность, точный двусторонний биномиальный тест Мак-Немара и скорректированное по Холму $p$ для семейства четырёх тестов при $\alpha=0.05$. Только эти четыре теста Мак-Немара определяют основное семейство проверок; бутстрэп-интервалы и тесты смены знаков не служат критерием основного отклонения гипотезы. Процентильные интервалы используют 10\,000 бутстрэп-выборок задач с начальным состоянием анализа 20260908; эти описательные интервалы не являются одновременными интервалами. Взаимодействие, ресурсные контрасты и сравнения с прямым решателем остаются описательными. Взаимодействие использует задачи с полными наблюдениями во всех четырёх факторных условиях; оно может не совпадать с вычитанием двух отдельно оценённых ролевых контрастов, если их доступные подвыборки различаются. Повторы разработки усредняются внутри задачи в отдельном описательном анализе и не объединяются с основными тестами.

При 200 независимых оцениваемых задачах наихудшая нормальная 95\%-полуширина для одной доли приблизительно равна 6.9 процентного пункта. Это обоснование точности, а не расчёт 80\%-мощности для парного эффекта; неопределённость парной разности зависит от несогласованных исходов. Приемлемая потеря качества для конкретного применения не получила внешнего обоснования, поэтому граница неухудшения и решение о сохранении качества не используются. Порог, выбранный только под статистическую точность, не устанавливал бы практически приемлемую потерю. Качество условия оценивается по доступным наблюдениям и сопровождается границами пропусков по всем назначениям: недоступные исходы считаются сначала всеми неуспехами, затем всеми успехами. Эти границы не являются доверительными интервалами или поправкой при случайных пропусках.

Завершённый вызов CLI возвращает счётчики входа $I$, кешированного входа $C$ и выхода $O$. Кешированный вход входит в общий вход; отдельно доступное число токенов рассуждения входит в выход. Повторного суммирования нет. Опубликованные стандартные тарифы API составляют соответственно USD 0.75, 0.075 и 4.50 за миллион токенов~\cite{openaipricing2026}. Рассчитываются
\begin{equation}
V=\frac{0.75(I-C)+0.075C+4.50O}{10^6},\qquad
V_0=\frac{0.75I+4.50O}{10^6}\quad\mathrm{USD}.
\end{equation}
$V$ --- контрфактическая оценка измеренных токенов подписки по прейскуранту API, а не платёж или счёт; $V_0$ --- анализ чувствительности без скидки за кеш. Обе оценки исключают работу исследовательских ассистентов, разработку кода, сборку окружения и вычисления оценщика. Средние ресурсы программы суммируют её завершённые вызовы. Отдельный реестр включает доступные счётчики прерванных программ и явно считает вызовы без итогового учёта; неизвестное использование не заменяется нулём. Отношения ресурсов используют одну и ту же подвыборку полных пар в числителе и знаменателе. «Меньшая наблюдаемая стоимость» означает отношение парных средних оценок ниже единицы, то есть отрицательную среднюю парную разность. Это описательный критерий: подтверждающий денежный порог и совместное правило решения о качестве и стоимости заранее не задавались; ресурсные интервалы не устанавливают денежное превосходство. Время отправки и состояние кеша могут влиять на денежную оценку, особенно между отдельно запущенными сериями.
